% DO NOT EDIT - automatically generated from metadata.yaml

\def \codeURL{https://github.com/VolMax-Studio/rotation-gap-fano-s1}
\def \codeDOI{10.5281/zenodo.22309778}
\def \codeSWH{}
\def \dataURL{https://doi.org/10.5281/zenodo.13273331}
\def \dataDOI{10.5281/zenodo.13273331}
\def \editorNAME{}
\def \editorORCID{}
\def \reviewerINAME{}
\def \reviewerIORCID{}
\def \reviewerIINAME{}
\def \reviewerIIORCID{}
\def \dateRECEIVED{}
\def \dateACCEPTED{}
\def \datePUBLISHED{}
\def \articleTITLE{[Re] Partial Replication of the Willow Fano-Factor Results in "The Rotation Gap Is Not An Error"}
\def \articleTYPE{Replication}
\def \articleDOMAIN{Quantum Computation and Statistical Physics}
\def \articleBIBLIOGRAPHY{bibliography.bib}
\def \articleYEAR{2026}
\def \reviewURL{}
\def \articleABSTRACT{We present a pre-registered partial replication of the empirical quantum error-correction
detector statistics reported in Section II F and Table III of Stenberg (2026, arXiv:2604.11963v1).
Analyzing raw detector bitstream records from the 105-qubit Google Willow archive (Zenodo DOI:
10.5281/zenodo.13273331, 420 experiments across code distances d in {3, 5, 7}), we recompute
all four primary and secondary claims under strict cryptographic and algorithmic verification gates.
We confirm that: (1) the aggregate empirical Fano factor reproduces to F = 2.4157 ≈ 2.42 ± 0.36 (Claim W1);
(2) distance-stratified Fano factors reproduce to d=3 -> 2.2942, d=5 -> 2.5935, and d=7 -> 2.7978 (Claim W2);
(3) the one-way ANOVA across distances yields F = 59.13, p = 2.46e-23 (Claim W3); and (4) the one-sample
t-statistic against Poisson expectation (F = 1) reproduces to t = +80.15 (Claim W4).
We explicitly document our three-layer mixed implementation provenance: the primary detector decoding path
is independently implemented (read_detection_counts_fast), while the statistical estimator on the analytical path
deliberately transcribes the author's published operations under the preregistered design, with reference
author decoding operations retained off-path for 10/10 exact bitwise equivalence verification.
Finally, we document two epistemic limitations: a structural zero-power defect in the preregistered
QEC round-sensitivity test resulting from Google's balanced experimental design (N_r = 28), and an exploratory
finding demonstrating that equal-distance weighting shifts the recomputed mean Fano factor from 2.4157 to 2.5618
(|Delta| = 0.1461, with the corresponding source-reported aggregate being 2.42) due to the predominant representation
of d=3 experiments (64.3%) in the public repository.
}
\def \replicationCITE{Stenberg, S. (2026). The Rotation Gap Is Not An Error — Ternary Structure in IBM Quantum Hardware. arXiv:2604.11963v1 [quant-ph].}
\def \replicationBIB{stenberg2026rotation}
\def \replicationURL{https://arxiv.org/abs/2604.11963}
\def \replicationDOI{10.48550/arXiv.2604.11963}
\def \contactNAME{Ivan Nestorov}
\def \contactEMAIL{volmax.core@gmail.com}
\def \articleKEYWORDS{quantum error correction, surface codes, google willow, fano factor, replication, python}
\def \journalNAME{ReScience C}
\def \journalVOLUME{}
\def \journalISSUE{}
\def \articleNUMBER{}
\def \articleDOI{}
\def \authorsFULL{Ivan Nestorov}
\def \authorsABBRV{I. Nestorov}
\def \authorsSHORT{Nestorov}
\title{\articleTITLE}
\date{}
\author[1,\orcid{0009-0006-7940-9539}]{Ivan Nestorov}
\affil[1]{VolMax Studio Lab, Belgrade, Serbia}
